\documentclass[10pt,a4paper]{amsart}
\usepackage{amsmath,amssymb,amsthm,graphicx}
\usepackage{verbatim}
\newenvironment{reference}{\comment}{\endcomment}
%\newenvironment{reference}{}{}
\newenvironment{slogan}{\comment}{\endcomment}
\newenvironment{history}{\comment}{\endcomment}
% For commutative diagrams we use Xy-pic
\usepackage[all]{xy}

% We use 2cell for 2-commutative diagrams.
\xyoption{2cell}
% Theorem environments.
%
\theoremstyle{plain}
\newtheorem{theorem}[subsection]{Theorem}
\newtheorem{proposition}[subsection]{Proposition}
\newtheorem{lemma}[subsection]{Lemma}

\theoremstyle{definition}
\newtheorem{definition}[subsection]{Definition}
\newtheorem{example}[subsection]{Example}
\newtheorem{exercise}[subsection]{Exercise}
\newtheorem{situation}[subsection]{Situation}

\theoremstyle{remark}
\newtheorem{remark}[subsection]{Remark}
\newtheorem{remarks}[subsection]{Remarks}

\numberwithin{equation}{subsection}

% Macros
%
\def\lim{\mathop{\mathrm{lim}}\nolimits}
\def\colim{\mathop{\mathrm{colim}}\nolimits}
\def\Spec{\mathop{\mathrm{Spec}}}
\def\Hom{\mathop{\mathrm{Hom}}\nolimits}
\def\Ext{\mathop{\mathrm{Ext}}\nolimits}
\def\SheafHom{\mathop{\mathcal{H}\!\mathit{om}}\nolimits}
\def\SheafExt{\mathop{\mathcal{E}\!\mathit{xt}}\nolimits}
\def\Sch{\mathit{Sch}}
\def\Mor{\mathop{\mathrm{Mor}}\nolimits}
\def\Ob{\mathop{\mathrm{Ob}}\nolimits}
\def\Sh{\mathop{\mathit{Sh}}\nolimits}
\def\NL{\mathop{N\!L}\nolimits}
\def\CH{\mathop{\mathrm{CH}}\nolimits}
\def\proetale{{pro\text{-}\acute{e}tale}}
\def\etale{{\acute{e}tale}}
\def\QCoh{\mathit{QCoh}}
\def\Ker{\mathop{\mathrm{Ker}}}
\def\Im{\mathop{\mathrm{Im}}}
\def\Coker{\mathop{\mathrm{Coker}}}
\def\Coim{\mathop{\mathrm{Coim}}}

% Boxtimes
%
\DeclareMathSymbol{\boxtimes}{\mathbin}{AMSa}{"02}

%
% Macros for moduli stacks/spaces
%
\def\QCohstack{\mathcal{QC}\!\mathit{oh}}
\def\Cohstack{\mathcal{C}\!\mathit{oh}}
\def\Spacesstack{\mathcal{S}\!\mathit{paces}}
\def\Quotfunctor{\mathrm{Quot}}
\def\Hilbfunctor{\mathrm{Hilb}}
\def\Curvesstack{\mathcal{C}\!\mathit{urves}}
\def\Polarizedstack{\mathcal{P}\!\mathit{olarized}}
\def\Complexesstack{\mathcal{C}\!\mathit{omplexes}}
% \Pic is the operator that assigns to X its picard group, usage \Pic(X)
% \Picardstack_{X/B} denotes the Picard stack of X over B
% \Picardfunctor_{X/B} denotes the Picard functor of X over B
\def\Pic{\mathop{\mathrm{Pic}}\nolimits}
\def\Picardstack{\mathcal{P}\!\mathit{ic}}
\def\Picardfunctor{\mathrm{Pic}}
\def\Deformationcategory{\mathcal{D}\!\mathit{ef}}
\usepackage{cleveref}
\crefname{originalreference}{original reference}{original references}
\Crefname{originalreference}{Original reference}{Original references}
\makeatletter
\newlabel{algebraic-definition-presentation}{{algebraic-definition-presentation}{source}{}{sourcepointer-1}{}}
\newlabel{algebraic-definition-presentation@cref}{{[originalreference][0][]algebraic-definition-presentation}{[1][0][]source}{}{}{}}
\makeatother
\crefname{originalreference}{original reference}{original references}
\Crefname{originalreference}{Original reference}{Original references}
\makeatletter
\newlabel{algebraic-lemma-map-space-into-stack}{{algebraic-lemma-map-space-into-stack}{source}{}{sourcepointer-2}{}}
\newlabel{algebraic-lemma-map-space-into-stack@cref}{{[originalreference][0][]algebraic-lemma-map-space-into-stack}{[1][0][]source}{}{}{}}
\makeatother
\crefname{originalreference}{original reference}{original references}
\Crefname{originalreference}{Original reference}{Original references}
\makeatletter
\newlabel{algebraic-lemma-stack-presentation}{{algebraic-lemma-stack-presentation}{source}{}{sourcepointer-3}{}}
\newlabel{algebraic-lemma-stack-presentation@cref}{{[originalreference][0][]algebraic-lemma-stack-presentation}{[1][0][]source}{}{}{}}
\makeatother
\crefname{originalreference}{original reference}{original references}
\Crefname{originalreference}{Original reference}{Original references}
\makeatletter
\newlabel{algebraic-theorem-smooth-groupoid-gives-algebraic-stack}{{algebraic-theorem-smooth-groupoid-gives-algebraic-stack}{source}{}{sourcepointer-4}{}}
\newlabel{algebraic-theorem-smooth-groupoid-gives-algebraic-stack@cref}{{[originalreference][0][]algebraic-theorem-smooth-groupoid-gives-algebraic-stack}{[1][0][]source}{}{}{}}
\makeatother
\crefname{originalreference}{original reference}{original references}
\Crefname{originalreference}{Original reference}{Original references}
\makeatletter
\newlabel{groupoids-section-groupoids}{{groupoids-section-groupoids}{source}{}{sourcepointer-5}{}}
\newlabel{groupoids-section-groupoids@cref}{{[originalreference][0][]groupoids-section-groupoids}{[1][0][]source}{}{}{}}
\makeatother
\crefname{originallemma}{original lemma}{original lemmas}
\Crefname{originallemma}{Original lemma}{Original lemmas}
\makeatletter
\newlabel{lemma-prepare-normalization}{{lemma-prepare-normalization}{source}{}{sourcepointer-6}{}}
\newlabel{lemma-prepare-normalization@cref}{{[originallemma][0][]lemma-prepare-normalization}{[1][0][]source}{}{}{}}
\makeatother
\crefname{originalreference}{original reference}{original references}
\Crefname{originalreference}{Original reference}{Original references}
\makeatletter
\newlabel{morphisms-definition-normalization}{{morphisms-definition-normalization}{source}{}{sourcepointer-7}{}}
\newlabel{morphisms-definition-normalization@cref}{{[originalreference][0][]morphisms-definition-normalization}{[1][0][]source}{}{}{}}
\makeatother
\crefname{originalreference}{original reference}{original references}
\Crefname{originalreference}{Original reference}{Original references}
\makeatletter
\newlabel{spaces-groupoids-lemma-criterion-fibre-product}{{spaces-groupoids-lemma-criterion-fibre-product}{source}{}{sourcepointer-8}{}}
\newlabel{spaces-groupoids-lemma-criterion-fibre-product@cref}{{[originalreference][0][]spaces-groupoids-lemma-criterion-fibre-product}{[1][0][]source}{}{}{}}
\makeatother
\crefname{originalreference}{original reference}{original references}
\Crefname{originalreference}{Original reference}{Original references}
\makeatletter
\newlabel{spaces-morphisms-lemma-normalization}{{spaces-morphisms-lemma-normalization}{source}{}{sourcepointer-9}{}}
\newlabel{spaces-morphisms-lemma-normalization@cref}{{[originalreference][0][]spaces-morphisms-lemma-normalization}{[1][0][]source}{}{}{}}
\makeatother
\title{A stack satisfying the normalization hypotheses has an integral normalization whose formation commutes with smooth base change}
\author{The Stacks Project Authors}
\date{stacks-project:stacks-morphisms.tex:local-snapshot-20261012; GNU FDL 1.2 or later}
\begin{document}
\maketitle
\noindent\textit{Editorial note.} All delivered mathematical text is the author's own, copied from the pinned Stacks Project chapter source sources/stacks-algebraic-stacks-morphisms/pinned-original-source.tex; no statement or proof step was rewritten, reordered or removed.\par\medskip
\noindent\textit{Editorial note.} Preamble provenance: the Stacks Project keeps one single shared preamble for every chapter; this card takes it directly from the staged sources/stacks-algebraic-stacks-morphisms/preamble.tex as three verbatim spans (lines 8-12, 14-18 and 154-217), so no derived file was written and nothing was added or redefined. The wrapper class and the base packages are supplied by the assembler.\par\medskip
\noindent\textit{Editorial note.} Bibliography provenance: neither the counted statement nor its proof contains a citation (this lane rescans the whole span before writing the card), so the card carries no bibliography and no entry was reconstructed.\par\medskip
\noindent\textit{Editorial note.} Selection provenance: this card was opened under the file-based control contract (CONTROL.md), with the mechanical pre-card gate run on its statement and proof interval before the card was written.\par\medskip
\noindent\textit{Allowed standard prerequisites (editorial).}\begin{itemize}
\item The preparation lemma describing the hypotheses on the stack in terms of a presentation. Invoked in the statement and at the start of the counted proof (source lines 10235-10249). Provides the presentation carrying the hypotheses that make normalization available.
\item The correspondence between algebraic stacks presented by groupoids and their quotients. Invoked inside the counted proof (source lines 10252-10270) as a cross-chapter pointer. Presents the stack as a quotient of a scheme by a groupoid so that normalizations may be glued.
\item The construction of the normalization of a scheme and its behaviour under smooth morphisms. Invoked inside the counted proof (source lines 10271-10295) as cross-chapter pointers. Produces the integral morphism from the normalized presentation.
\end{itemize}\medskip
\ifdefined\hypertarget\hypertarget{sourcepointer-1}{}\fi
\noindent\textit{Source reference.} Original reference algebraic-definition-presentation: sources/stacks-algebraic-stacks-morphisms/pinned-original-source.tex, source lines 10255--10255. Navigation pointer only: the counted text refers to the Stacks item algebraic-definition-presentation, which belongs to another chapter of the project and is not part of this card. The pointer carries the target's own label as its number because this chapter's staged snapshot does not contain the other chapter's tag table; its content is not reproduced here.\par\medskip
\ifdefined\hypertarget\hypertarget{sourcepointer-2}{}\fi
\noindent\textit{Source reference.} Original reference algebraic-lemma-map-space-into-stack: sources/stacks-algebraic-stacks-morphisms/pinned-original-source.tex, source lines 3831--3831. Navigation pointer only: the counted text refers to the Stacks item algebraic-lemma-map-space-into-stack, which belongs to another chapter of the project and is not part of this card. The pointer carries the target's own label as its number because this chapter's staged snapshot does not contain the other chapter's tag table; its content is not reproduced here.\par\medskip
\ifdefined\hypertarget\hypertarget{sourcepointer-3}{}\fi
\noindent\textit{Source reference.} Original reference algebraic-lemma-stack-presentation: sources/stacks-algebraic-stacks-morphisms/pinned-original-source.tex, source lines 287--287. Navigation pointer only: the counted text refers to the Stacks item algebraic-lemma-stack-presentation, which belongs to another chapter of the project and is not part of this card. The pointer carries the target's own label as its number because this chapter's staged snapshot does not contain the other chapter's tag table; its content is not reproduced here.\par\medskip
\ifdefined\hypertarget\hypertarget{sourcepointer-4}{}\fi
\noindent\textit{Source reference.} Original reference algebraic-theorem-smooth-groupoid-gives-algebraic-stack: sources/stacks-algebraic-stacks-morphisms/pinned-original-source.tex, source lines 10291--10291. Navigation pointer only: the counted text refers to the Stacks item algebraic-theorem-smooth-groupoid-gives-algebraic-stack, which belongs to another chapter of the project and is not part of this card. The pointer carries the target's own label as its number because this chapter's staged snapshot does not contain the other chapter's tag table; its content is not reproduced here.\par\medskip
\ifdefined\hypertarget\hypertarget{sourcepointer-5}{}\fi
\noindent\textit{Source reference.} Original reference groupoids-section-groupoids: sources/stacks-algebraic-stacks-morphisms/pinned-original-source.tex, source lines 10277--10277. Navigation pointer only: the counted text refers to the Stacks item groupoids-section-groupoids, which belongs to another chapter of the project and is not part of this card. The pointer carries the target's own label as its number because this chapter's staged snapshot does not contain the other chapter's tag table; its content is not reproduced here.\par\medskip
\ifdefined\hypertarget\hypertarget{sourcepointer-6}{}\fi
\noindent\textit{Source reference.} Original lemma lemma-prepare-normalization: sources/stacks-algebraic-stacks-morphisms/pinned-original-source.tex, source lines 10198--10198. Navigation pointer only: the counted text refers to the Stacks item lemma-prepare-normalization, which is an existing labeled result of this same chapter and is not part of this card. The pointer carries the target's own label as its number and its definition line as its span; its content is not reproduced here.\par\medskip
\ifdefined\hypertarget\hypertarget{sourcepointer-7}{}\fi
\noindent\textit{Source reference.} Original reference morphisms-definition-normalization: sources/stacks-algebraic-stacks-morphisms/pinned-original-source.tex, source lines 10258--10258. Navigation pointer only: the counted text refers to the Stacks item morphisms-definition-normalization, which belongs to another chapter of the project and is not part of this card. The pointer carries the target's own label as its number because this chapter's staged snapshot does not contain the other chapter's tag table; its content is not reproduced here.\par\medskip
\ifdefined\hypertarget\hypertarget{sourcepointer-8}{}\fi
\noindent\textit{Source reference.} Original reference spaces-groupoids-lemma-criterion-fibre-product: sources/stacks-algebraic-stacks-morphisms/pinned-original-source.tex, source lines 10299--10299. Navigation pointer only: the counted text refers to the Stacks item spaces-groupoids-lemma-criterion-fibre-product, which belongs to another chapter of the project and is not part of this card. The pointer carries the target's own label as its number because this chapter's staged snapshot does not contain the other chapter's tag table; its content is not reproduced here.\par\medskip
\ifdefined\hypertarget\hypertarget{sourcepointer-9}{}\fi
\noindent\textit{Source reference.} Original reference spaces-morphisms-lemma-normalization: sources/stacks-algebraic-stacks-morphisms/pinned-original-source.tex, source lines 10260--10260. Navigation pointer only: the counted text refers to the Stacks item spaces-morphisms-lemma-normalization, which belongs to another chapter of the project and is not part of this card. The pointer carries the target's own label as its number because this chapter's staged snapshot does not contain the other chapter's tag table; its content is not reproduced here.\par\medskip
\begin{lemma}
\label{lemma-normalization}
Let $\mathcal{X}$ be an algebraic stack satisfying
the equivalent conditions of Lemma \ref{lemma-prepare-normalization}.
Then there exists an integral morphism of algebraic stacks
$$
\mathcal{X}^\nu \longrightarrow \mathcal{X}
$$
such that for every scheme $U$ and smooth morphism $U \to \mathcal{X}$
the fibre product $\mathcal{X}^\nu \times_\mathcal{X} U$
is the normalization of $U$.
\end{lemma}

\begin{proof}
Let $U \to \mathcal{X}$ be a surjective smooth morphism where $U$ is a scheme.
Set $R = U \times_\mathcal{X} U$. Recall that we obtain a smooth groupoid
$(U, R, s, t, c)$ in algebraic spaces and a presentation
$\mathcal{X} = [U/R]$ of $\mathcal{X}$, see Algebraic Stacks,
Lemmas \ref{algebraic-lemma-map-space-into-stack} and
\ref{algebraic-lemma-stack-presentation} and
Definition \ref{algebraic-definition-presentation}.
The assumption on $\mathcal{X}$ means that the normalization
$U^\nu$ of $U$ is defined, see Morphisms, Definition
\ref{morphisms-definition-normalization}.
By Morphisms of Spaces, Lemma
\ref{spaces-morphisms-lemma-normalization}
taking normalization commutes with smooth morphisms of algebraic spaces.
Thus we see that the normalization $R^\nu$ of $R$
is isomorphic to both $R \times_{s, U} U^\nu$ and $U^\nu \times_{U, t} R$.
Thus we obtain two smooth morphisms
$s^\nu : R^\nu \to U^\nu$ and $t^\nu : R^\nu \to U^\nu$
of algebraic spaces. A formal computation with fibre products shows that
$R^\nu \times_{s^\nu, U^\nu, t^\nu} R^\nu$ is the normalization
of $R \times_{s, U, t} R$. Hence the smooth morphism
$c : R \times_{s, U, t} R \to R$ extends to $c^\nu$ as well.
Similarly, the inverse $i : R \to R$ (an isomorphism) induces
an isomorphism $i^\nu : R^\nu \to R^\nu$. Finally, the identity
$e : U \to R$ lifts to $e^\nu : U^\nu \to R^\nu$ for example
because $e$ is a section of $s$ and $R^\nu = R \times_{U, s} U^\nu$.
We claim that $(U^\nu, R^\nu, s^\nu, t^\nu, c^\nu)$ is
a smooth groupoid in algebraic spaces. To see this involves checking
the axioms (1), (2)(a), (2)(b), (3)(a), and (3)(b) of
Groupoids, Section \ref{groupoids-section-groupoids}
for $(U^\nu, R^\nu, s^\nu, t^\nu, c^\nu, e^\nu, i^\nu)$.
For example, for (1) we have to see that the two morphisms
$a, b : R^\nu \times_{s^\nu, U^\nu, t^\nu} R^\nu
\times_{s^\nu, U^\nu, t^\nu} R^\nu \to R^\nu$ we obtain are the same.
This holds because we know that the corresponding pair of morphisms
$R \times_{s, U, t} R \times_{s, U, t} R \to R$
are the same and the morphisms $a$ and $b$ are the unique extensions
of this morphism to the normalizations. Similarly for the other
axioms.

\medskip\noindent
Consider the algebraic stack $\mathcal{X}^\nu = [U^\nu/R^\nu]$
(Algebraic Stacks, Theorem
\ref{algebraic-theorem-smooth-groupoid-gives-algebraic-stack}).
Since we have a morphism
$(U^\nu, R^\nu, s^\nu, t^\nu, c^\nu) \to (U, R, s, t, c)$
of groupoids in algebraic spaces, we obtain a morphism
$\nu : \mathcal{X}^\nu \to \mathcal{X}$ of algebraic stacks.
Since $R^\nu = R \times_{s, U} U^\nu$ we see that
$U^\nu = \mathcal{X}^\nu \times_\mathcal{X} U$ by
Groupoids in Spaces, Lemma
\ref{spaces-groupoids-lemma-criterion-fibre-product}.
In particular, as $U^\nu \to U$ is integral, we see that $\nu$ is integral.
We omit the verification that the base change property
stated in the lemma holds for every smooth morphism from a scheme to
$\mathcal{X}$.
\end{proof}

\end{document}
